\chapter[Plan de Capacitación y Entrenamiento]{Plan de Capacitación y Entrenamiento}

La capacitación se dirige al equipo de Enfermería de la Universidad del Bío-Bío,
que recibe el prototipo al cierre del proyecto y lo presentará luego a los
cuidadores. Capacitar primero a este equipo permite que las familias aprendan a
usar la aplicación de la mano de profesionales que ya conocen su contenido
clínico.

\begin{itemize}
    \item \emph{Usuarios a capacitar}: académicas y tesistas de Enfermería que
    acompañarán la puesta en marcha.
    \item \emph{Tipo de capacitación}: taller práctico presencial en que cada
    participante usa la aplicación en su propio teléfono, apoyado por el manual
    de usuario del anexo.
    \item \emph{Funcionalidades abordadas}: creación de cuenta y pacientes,
    registro clínico y niveles de alerta, historial y exportación, biblioteca,
    asistente, preguntas frecuentes y recordatorios.
    \item \emph{Responsables}: las estudiantes desarrolladoras.
    \item \emph{Tiempo estimado}: dos sesiones de una hora y media.
    \item \emph{Recursos}: teléfonos con la aplicación instalada, proyector,
    manual de usuario y una cuenta de demostración con datos ficticios.
\end{itemize}

\begin{table}[H]
\centering
\small
\tablazebra
\caption{Programa de capacitación.}
\label{tab:capacitacion}
\begin{tabularx}{\textwidth}{P{0.55}P{0.75}P{2.05}P{0.65}}
\toprule
\filaencabezado \encab{Sesión} & \encab{Fecha} & \encab{Contenido} & \encab{Duración} \\
\midrule
1 & 16 nov 2026 & Instalación, creación de cuenta y pacientes, registro
clínico, niveles de alerta, historial y exportación. & 1,5 h \\
2 & 18 nov 2026 & Biblioteca, asistente, preguntas frecuentes, recordatorios y
resolución de dudas. & 1,5 h \\
\bottomrule
\end{tabularx}
\end{table}
